% apendices/patrones-confusos.tex
\chapter{\ruby{紛}{まぎ}らわしい\ruby{パターン}{パターン} — Patrones Confusos}

\section*{4\ruby{つ}{つ}の\ruby{条件形}{じょうけんけい} — Los 4 Condicionales}

\begin{tcolorbox}[comparacionbox, title={⚖ たら vs ば vs なら vs と}]
\begin{tabularx}{\linewidth}{lXXl}
\toprule
\textbf{Forma} & \textbf{Cuándo usarlo} & \textbf{No usar cuando} & \textbf{Ej. clave} \\ \midrule
たら & Secuencia futura, hipótesis, descubrimiento & — (el más flexible) & \ruby{着}{つ}いたら\ruby{電話}{でんわ}する \\[4pt]
ば & Condición general, proverbio, 〜ばよかった & Resultado volitivo 1ª persona & \ruby{急}{いそ}がば\ruby{回}{まわ}れ \\[4pt]
なら & Responder a info del interlocutor, consejo & Evento propio futuro (usa たら) & \ruby{行}{い}くなら\ruby{傘}{かさ}を \\[4pt]
と & Resultado automático/inevitable, instrucciones & Resultado volitivo («voy a hacer X») & ボタンを\ruby{押}{お}すと\ruby{光}{ひか}る \\
\bottomrule
\end{tabularx}
\end{tcolorbox}

\section*{\ruby{証拠}{しょうこ}・\ruby{推量}{すいりょう} — Evidencialidad}

\begin{tcolorbox}[comparacionbox, title={⚖ そうだ vs らしい vs ようだ vs はずだ}]
\begin{tabularx}{\linewidth}{lXX}
\toprule
\textbf{Forma} & \textbf{Fuente de info} & \textbf{Ejemplo} \\ \midrule
V-stem + そうだ & Apariencia visual directa & \ruby{雨}{あめ}が\ruby{降}{ふ}りそうだ \\
plain + そうだ & Información de otros (reportar) & \ruby{雨}{あめ}が\ruby{降}{ふ}るそうだ \\
らしい & Evidencia / rumores / impresión & \ruby{雨}{あめ}が\ruby{降}{ふ}るらしい \\
ようだ & Observación / inferencia directa & \ruby{雨}{あめ}が\ruby{降}{ふ}るようだ \\
はずだ & Expectativa lógica / deber ser & \ruby{今}{いま}\ruby{雨}{あめ}のはずだ \\
\bottomrule
\end{tabularx}
\end{tcolorbox}

\section*{\ruby{限定}{げんてい}と\ruby{否定}{ひてい} — Limitadores}

\begin{tcolorbox}[comparacionbox, title={⚖ だけ vs しか vs ばかり}]
\begin{tabularx}{\linewidth}{lXXl}
\toprule
\textbf{Forma} & \textbf{Significado} & \textbf{Frase verbal} & \textbf{Énfasis} \\ \midrule
だけ & solo / únicamente (neutro) & + afirmativo & limitación \\
しか & solo (excluye todo lo demás) & + nega\-tivo & exclusión total \\
ばかり & solo / nada más que (subjetivo) & + afirmativo & demasiado \\
\bottomrule
\end{tabularx}
\tcblower
\small \ruby{一人}{ひとり}だけ\ruby{来}{き}た。(Solo vino uno — neutro) \quad
\ruby{一人}{ひとり}しか\ruby{来}{こ}なかった。(Solo vino uno — decepcionante) \quad
\ruby{一人}{ひとり}ばかり\ruby{話}{はな}す。(Habla él solo todo el tiempo — queja)
\end{tcolorbox}

\section*{\ruby{意味}{いみ}が\ruby{似}{に}ている\ruby{表現}{ひょうげん} — Expresiones de Significado Similar}

\begin{tcolorbox}[comparacionbox, title={⚖ ために vs ように}]
\begin{tabularx}{\linewidth}{lXX}
\toprule
\textbf{Forma} & \textbf{Propósito/Meta} & \textbf{Sujeto de la meta} \\ \midrule
ために & meta concreta, acción deliberada & mismo sujeto / objeto directo \\
ように & estado deseado, cambio gradual & puede ser otro \\
\bottomrule
\end{tabularx}
\tcblower
\small \ruby{日本語}{にほんご}を\ruby{話}{はな}せるようにする。(\ruby{能力}{のうりょく}の\ruby{変化}{へんか}) \quad
\ruby{日本語}{にほんご}を\ruby{話}{はな}すために\ruby{勉強}{べんきょう}する。(\ruby{具体的}{ぐたいてき}\ruby{目的}{もくてき})
\end{tcolorbox}

\begin{tcolorbox}[comparacionbox, title={⚖ こと vs の (nominalización)}]
\begin{tabularx}{\linewidth}{lXX}
\toprule
\textbf{Forma} & \textbf{Uso preferido} & \textbf{Ejemplo} \\ \midrule
こと & concepto abstracto, regla, formal & \ruby{運動}{うんどう}することが\ruby{大切}{たいせつ}だ \\
の & situación concreta, percepción & \ruby{彼}{かれ}が\ruby{来}{く}るのが\ruby{見}{み}える \\
\bottomrule
\end{tabularx}
\end{tcolorbox}

\section*{\ruby{て形}{てけい}の\ruby{使役受身}{しえきうけみ} — Cadena Pasiva-Causativa}

\begin{tcolorbox}[vocabbox, title={\ruby{原形}{げんけい} → \ruby{使役形}{しえきけい} → \ruby{受身形}{うけみけい} → \ruby{使役受身形}{しえきうけみけい}}]
\begin{tabularx}{\linewidth}{lXXX}
\toprule
\textbf{G1 ejemplo} & \textbf{飲む} & \textbf{書く} & \textbf{待つ} \\ \midrule
辞書形 & \ruby{飲}{の}む & \ruby{書}{か}く & \ruby{待}{ま}つ \\
使役形 & \ruby{飲}{の}ませる & \ruby{書}{か}かせる & \ruby{待}{ま}たせる \\
受身形 & \ruby{飲}{の}まれる & \ruby{書}{か}かれる & \ruby{待}{ま}たれる \\
使役受身 & \ruby{飲}{の}まされる & \ruby{書}{か}かされる & \ruby{待}{ま}たされる \\
\midrule
\textbf{G2 ejemplo} & \textbf{食べる} & \textbf{見る} & \textbf{教える} \\ \midrule
辞書形 & \ruby{食}{た}べる & \ruby{見}{み}る & \ruby{教}{おし}える \\
使役形 & \ruby{食}{た}べさせる & \ruby{見}{み}させる & \ruby{教}{おし}えさせる \\
受身形 & \ruby{食}{た}べられる & \ruby{見}{み}られる & \ruby{教}{おし}えられる \\
使役受身 & \ruby{食}{た}べさせられる & \ruby{見}{み}させられる & \ruby{教}{おし}えさせられる \\
\bottomrule
\end{tabularx}
\end{tcolorbox}
